
\section{Teoremas de existencia, unicidad y compatibilidad de las constantes genealógicas}

La construcción determina un sector aut\'onomo de constantes intr\'insecas.
Sus composiciones, teoremas cruzados y certificados forman una cadena
probatoria autocontenida. El orden rector es

{\small
\begin{equation}
\APP\longrightarrow\TRIT\longrightarrow\TPK
\longrightarrow\text{realizaci\'on funcional}
\longrightarrow\text{invariante}
\longrightarrow\text{secci\'on}
\longrightarrow\text{valor terminal}.
\label{eq:final-causal-chain}
\end{equation}
}

Con ese orden se obtienen cinco cierres principales.

\begin{enumerate}
\item El vac\'io electromagn\'etico deja de ser una colecci\'on de cuatro
f\'ormulas: es un operador positivo de dos hojas cuyo espectro determina
\(\varepsilon,\mu,Z,c\), con amplificaci\'on non\'adica compatible.

\item La criticidad deja de comenzar en los decimales de Feigenbaum: la
dodecafase cancela \(\pi\), el modo \(30\) produce \(899\), el germen se cierra
trigonom\'etricamente y entra de forma demostrada en la clase anal\'itica
unimodal apropiada. La hiperbolicidad universal queda aislada como un
certificado distinto.

\item La gravedad deja de ser s\'olo una carta decimal: la periodicidad temporal de \(108\) fases selecciona una
secci\'on de \(L_G\), el proyector de rango cinco fija el portador tensorial y
la reducci\'on TT distingue ese portador de las dos helicidades sin masa.

\item El cambio hexada--octada adquiere tres realizaciones relacionadas:
cuanto
areal, memoria modular y reflexi\'on de mezcla. El Hamiltoniano hologr\'afico
modular es central en esa interfaz de confluencia, pero no absorbe el vac\'io, la criticidad,
la gravedad, la aritm\'etica ni la cosmolog\'ia en un t\'itulo artificialmente
\'unico.

\item Las constantes menos usuales incluyen Catalan, L\'evy, Lochs y
Euler--Kronecker. Bendersky--Glaisher, vacancias y las realizaciones cosmol\'ogicas
aparecen asimismo como valores de operadores y din\'amicas, no como una lista
a la que se buscan coincidencias retrospectivas.
\end{enumerate}

\section{Car\'acter teorem\'atico del corolario metrol\'ogico}

El \cref{cor:metrological-confluence} y la
\cref{tab:metrological-confluence} constituyen el cierre probatorio del
desarrollo. No son un resumen abreviado del material anterior. Para que una fila
entre en ese corolario debe conservar simult\'aneamente
\[
\text{antecedente}\mid\text{ecuaci\'on}\mid\text{invariante}
\mid\text{significado}\mid\text{valor terminal}.
\]
Si se elimina cualquiera de esas cinco componentes, la fila deja de ser una
salida HMT y se convierte en una mera equivalencia num\'erica. El estatuto y
el criterio de refutación no son columnas sustitutas: auditan desde fuera la fuerza y el
dominio de la qu\'intupla completa.

Esta formulaci\'on produce tres consecuencias estructurales y matemáticas.
Primero, el valor terminal puede ser una expresi\'on exacta en una recta de
magnitud; no se exige reducirla a SI para reconocerla. Segundo, un antecedente
puede admitir varias realizaciones funcionales ---por ejemplo, \(T\)
determina tanto la respuesta de vac\'io como \(D_A\)--- sin que dichas
realizaciones se identifiquen.
Tercero, un entero compartido como \(30\), una raz\'on compartida como \(4/3\)
o un escalar compartido como \(\log3\) s\'olo establece confluencia cuando se
exhibe el mapa que lo transporta y el invariante que conserva.

En particular, el atlas permite leer dos ecuaciones de cierre que antes
quedaban dispersas:
\begin{equation}
\boxed{E_P=k_B\Theta_{\rm clk}\log3},
\qquad
\boxed{
(6\to8)\longrightarrow
\bigl(\gammaH\ell_P^2,\,-\log\rho_A,\,R_{\rm CKM}\bigr).}
\label{eq:status-two-confluences}
\end{equation}
La primera hace confluir cuanto de energ\'ia, temperatura e informaci\'on. La
segunda determina, desde una interfaz de confluencia com\'un, un cuanto
geom\'etrico, un generador
modular y una reflexi\'on de mezcla. Ninguna ecuaci\'on afirma que sus
componentes sean el mismo objeto; afirma que sus genealog\'ias se han
conservado hasta el punto exacto de ramificaci\'on.

El control negativo tambi\'en es inmediato. Si \(E_P\) se usa para seleccionar
\(k_B\), si CKM se usa para recalibrar \(\gammaH\), si \(\log3\) se confunde
con una matriz densidad o si la igualdad de dos decimales sustituye al
entrelazador, el corolario no se aplica. La autonomía de la derivación consiste en
que estos controles, al igual que las derivaciones, est\'an formulados dentro
del propio desarrollo.

\section{Estratificación lógica de identidades internas, realizaciones dimensionales y reconocimientos externos}

La fuerza de una afirmaci\'on no se deduce de lo extraordinario de su
conclusi\'on ni de su proximidad con un valor conocido. Se deduce de la cadena
que la sostiene.

\begin{longtable}{p{36mm}p{96mm}}
\toprule
Estatuto & Contenido del desarrollo\\
\midrule
\emph{Teorema interno exacto}
& Operador de vac\'io y sus constitutivas; modo \(30\); carga y cocientes
cu\'anticos; escala t\'ermica; perfil cerrado del caos; ret\'icula
\(3/4/12\); identidades de Gauss; selector \(\theta_{108}\); familia de
Planck; reflexi\'on CKM; ley
\(M_n=(a_n/a_0)^{-6}\) en su carta.\\

\emph{Composici\'on matem\'atica exacta}
& Wien, Stefan--Boltzmann, gas fot\'onico, Euler--Kronecker,
Bendersky--Glaisher, relaci\'on angular de Weinberg, reconstrucci\'on conjunta
\'area--memoria y autoacoplamiento de Higgs en la interfaz de \'arbol.\\

\emph{Demostrado con estructura de partida expl\'icita}
& Resultados que dependen de una carta declarada, una orientaci\'on, una
medida, un estado de borde o un dato de frontera, sin introducir el valor
terminal como selector.\\

\emph{Teorema con hip\'otesis expl\'icita}
& Clasificaci\'on del cierre GNS como factor hiperinfinito
\(III_{\lambda_\partial}\), bajo fidelidad, compatibilidad y persistencia
infinita de los dos canales \(90/120\).\\

\emph{Realizaci\'on f\'isica condicionada}
& Identificaci\'on completa de calibre, \(G_F\) renormalizada, realizaci\'on de
Sitter del
radio CT108, transducci\'on de la memoria de Perron a densidad de esp\'in y
realizaci\'on din\'amica del portador tensorial.\\

\emph{Programa abierto con criterio de refutación}
& Hiperbolicidad completa del renormalizador de Feigenbaum, autovalor
subdominante Gauss--Kuzmin--Wirsing y construcci\'on de una acci\'on
Born--Infeld seleccionada internamente.\\

\emph{Contraste metrol\'ogico externo}
& Cartas SI, comparaciones CODATA/PDG, valores cl\'asicos de las constantes y
clasificaciones operatorias aplicadas s\'olo despu\'es de construir los pesos
HMT.\\
\bottomrule
\end{longtable}

Una interfaz abierta no rebaja el teorema interno que la antecede. A la
inversa, un teorema interno no promueve autom\'aticamente una interpretaci\'on
f\'isica cuyo transductor a\'un no se haya construido.

\section{Obstrucciones transversales de tipado, causalidad, orientación y compatibilidad nonádica}

La construcción queda sometida a refutaci\'on. Los siguientes
controles atraviesan todos los cap\'itulos.

\begin{enumerate}[label=F\arabic*.]
\item \textbf{Tipado dimensional.} Una identidad que no sea homog\'enea en las
rectas \(L_S,L_C,L_T,L_Q\) queda descartada, aunque reproduzca un decimal.

\item \textbf{No alimentaci\'on por el valor objetivo.} Ning\'un valor SI, CODATA, PDG,
Feigenbaum, GKW, \(H_0\), \(G_F\) o anomal\'ia magn\'etica puede seleccionar
un operador, coeficiente o realizaci\'on interna.

\item \textbf{Convenci\'on de hojas.} Invertir silenciosamente \(q_+\) y
\(q_-\), o confundir la involuci\'on de hoja con conjugaci\'on CP, invalida las
paridades y los cocientes.

\item \textbf{Compatibilidad non\'adica.} Un funcional que cambie bajo
amplificaci\'on \(9^m\), pierda la unidad o no conserve la memoria de retorno
no posee el l\'imite reclamado.

\item \textbf{No fusi\'on por escalar.} La igualdad de dos valores terminales
no identifica sus dominios. Hace falta un entrelazador que conmute con los
operadores y preserve orientaci\'on, acarreo y frontera.

\item \textbf{Separaci\'on finito--infinito.} Un certificado hasta nivel ocho
no prueba la hiperbolicidad universal; un Hamiltoniano modular finito no es
\(-\log\rho_\infty\) dentro de la UHF; una convergencia num\'erica sin cota de
cola no certifica un producto primo.

\item \textbf{Transductor f\'isico.} La lectura de Sitter, Einstein--Cartan,
Born--Infeld o de calibre queda bloqueada si no determina el objeto completo de la
teor\'ia de destino y sus condiciones de borde.
\end{enumerate}

\section{Condiciones necesarias para las extensiones matemáticas y realizaciones físicas no equivalentes}

No queda una cautela gen\'erica del tipo ``falta relacionar HMT con la
f\'isica''. Quedan obligaciones localizadas:

\begin{description}[style=nextline,leftmargin=37mm]
\item[CHAOS-OBL-1]
Certificar en un espacio de anal\'iticas pares el punto fijo de
\(\mathcal R\), un \'unico autovalor inestable, la cota de cola y la
transversalidad del germen HMT.

\item[GAUSS-OBL-1]
Fijar el espacio funcional del operador de Gauss, aislar el autovalor
subdominante con intervalos y demostrar estabilidad del proyector bajo las
particiones \(9^m\).

\item[HMOD-OBL-1]
Verificar la compatibilidad de los estados locales y la persistencia con
densidad positiva de ambos canales antes de aplicar la clasificaci\'on
Araki--Woods. No usar una densidad traza--clase global inexistente.

\item[EW-OBL-1]
Construir el entrelazador entre hojas HMT y base neutra/cargada electrod\'ebil,
y derivar \(\Delta r\) sin emplear el valor experimental de \(G_F\).

\item[COS-OBL-1]
Para de Sitter, demostrar que el realizador PCH selecciona din\'amicamente
\(R_{108}\). Para Einstein--Cartan, construir amplitud, corriente, signo y
carta dimensional del mapa \(M_n\mapsto s_n^2\).

\item[BORN-OBL-1]
Construir la acci\'on determinantal Born--Infeld y demostrar que la escala
\(E_*\) es su par\'ametro, en vez de identificarla por dimensiones.

\item[ANOM-OBL-1]
Cerrar la procedencia local de \(1000\) y \(54+1\) antes de promover la
evaluaci\'on de \(a_\mu\) a derivaci\'on completa.
\end{description}

Cada obligaci\'on posee un objeto, un dominio y un criterio de refutación. Ninguna autoriza a
reabrir los bloques exactos que la preceden.

\section{Delimitación formal del dominio de las derivaciones genealógicas}

\begin{itemize}
\item No rederiva \(\pi,e,\varphi,\alpha_{\HMT}\); los usa como entradas ya
construidas cuando una composici\'on posterior lo exige.
\item No compite con la Ley General de Masas ni recalcula la masa electr\'onica.
\item No rederiva \(\gammaH\); desarrolla su significado areal,
informacional y modular.
\item No identifica el portador de cinco componentes con cinco polarizaciones
de un gravit\'on sin masa.
\item No confunde una aproximaci\'on finita de Feigenbaum con el teorema
universal ni una escala dimensional con una acci\'on Born--Infeld.
\item No convierte el Hamiltoniano hologr\'afico modular en el nombre o la
explicaci\'on exclusiva de todas las constantes. Es una pieza central de la
interfaz de confluencia \'area--informaci\'on--mezcla dentro de un atlas m\'as amplio.
\end{itemize}

\section{Capacidad reconstructiva de la genealogía común sobre constantes matemáticas y físicas}

La singularidad metrol\'ogica de HMT--MD no consiste en producir muchos
decimales. Consiste en que el valor real aparece al final de una secuencia que
conserva memoria. La permeabilidad y la permitividad son los dos autovalores
cuadr\'aticos de una misma respuesta; \(G\) es la secci\'on que cierra un radio
Compton con un radio gravitatorio; \(\log3\) es simult\'aneamente contenido de
decisi\'on y raz\'on t\'ermica; Catalan es un momento del cuarto de torsi\'on;
\(\gammaH\) orienta el cuanto de \'area; el Hamiltoniano modular conserva qu\'e
canal produjo la ocupaci\'on; la reflexi\'on CKM conserva la coordenada
responsable del cambio de hoja; el
exponente \(-6\) conserva la memoria estable al determinar la escala
cosmol\'ogica.

\begin{proposition}[Resultado principal]
La significaci\'on estructural demostrada en este desarrollo se expresa
mediante una regla de realizaci\'on: un mismo antecedente puede admitir
varios funcionales sin perder su genealog\'ia, y cada uno determina
simult\'aneamente el significado de la constante y la raz\'on de su valor
terminal.
\end{proposition}

La afirmaci\'on final es, por ello, m\'as precisa que ``HMT reproduce
constantes''. HMT--MD construye operadores, caracteres, secciones y
transductores de los que las constantes son evaluaciones. Donde la cadena
queda cerrada, el valor y su significado se demuestran juntos. Donde queda una
interfaz, el criterio de refutación permite continuar sin degradar lo ya establecido ni
presentar como teorema lo que a\'un es programa.
